SCON과 SCCC를 바라보던 문성이는 흥미로운 규칙을 발견했다! 바로 두 문자열에서 문자 `S'와 `C'의 개수를 합하면 항상 짝수가 된다는 것이다. 이를 본 문성이는 문득 알파벳 대문자로 구성된 길이가 $N$인 문자열 중에서 `S'와 `C'의 개수의 합이 짝수인 문자열이 몇 개나 될 지 궁금해졌다.

문성이의 궁금증을 해결해 주자!